% !PRINTED 270
\weberpage{270}
Wir setzen
\weq{3}{\begin{aligned}
2A &= uv + (-1)^{\frac{n+1}{8}}(u_1v_1+u_2v_2),\\
B &= uvu_1v_1 + uvu_2v_2 + (-1)^{\frac{n+1}{8}}u_1v_1u_2v_2\\
&= \frac{2}{u_1v_1} + \frac{2}{u_2v_2} + (-1)^{\frac{n+1}{8}}\frac{2}{uv},
\end{aligned}}
so da\ss{} sich die zusammengeh\"origen Vertauschungen ergeben
\weq{4}{\begin{array}{lll}
\omega, & A, & B,\\
-\dfrac{1}{\omega}, & \varrho\,A, & \varrho^2 B,\\
\omega+1, & e^{\frac{\pi i(n+1)}{12}}\sigma\,A, & e^{-\frac{\pi i(n+1)}{12}}\sigma^2 B.
\end{array}}

Ein Produkt von der Form $A^h B^k$ nimmt also durch die beiden Vertauschungen
\[
\Bigl(\omega, -\frac{1}{\omega}\Bigr), \quad (\omega, \omega+1)
\]
die Faktoren an
\[
\varrho^{h-k}, \quad e^{(h-k)\frac{(n+1)\pi i}{12}}\sigma^{h-k},
\]
die $=1$ sind, wenn $n+1$ durch 3 teilbar ist. Wir bilden also jetzt mit numerischen Koeffizienten $M_{h,k}$ Funktionen der Form
\weq{5}{\Phi_{a,c,\partial} = \Sigma\,M_{h,k}A^h B^k,}
worin, wenn $n \equiv 0$ oder $\equiv 1$ (mod 3) ist, $h$ und $k$ nur solche (ganzzahlige) Werte annehmen d\"urfen, deren Differenzen $h-k$ bei der Teilung mit 3 denselben Rest lassen, wenn aber $n \equiv -1$ (mod 3) dieser Beschr\"ankung nicht unterworfen sind. Eine solche Funktion selbst, oder wenigstens ihre dritte Potenz gen\"ugt also einer Transformationsgleichung erster Stufe. Sie bleibt au\ss{}erdem f\"ur alle endlichen Werte von $j(\omega)$ endlich und ist folglich eine \gesperrt{ganze} algebraische Funktion von $j(\omega)$. [Die Transformation zweiter Ordnung $\begin{pmatrix} 1, & 1\\ -1, & 1 \end{pmatrix}$ braucht hier nicht zugezogen zu werden, da der Inbegriff der $\Phi_{a,c,\partial}$ schon bei der Substitution $(\omega, \omega+1)$, nicht erst bei $(\omega, \omega+2)$ unge\"andert bleibt. Vgl. die Bemerkung am Anfang des §\,73.]

Wenn wir also die Konstanten in $\Phi_{a,c,\partial}$ so bestimmen, da\ss{} in den Entwickelungen dieser Funktionen keine negativen Potenzen von $q$ vorkommen, so m\"ussen die s\"amtlichen $\Phi_{a,c,\partial}$ einer und derselben Konstanten gleich sein.
